\section{问题四：开放模型能力与算力放缓情景}
\subsection{样本定义和三个可比性边界}
采用 C1 六任务同协议均分 $S$，按 C2 中权重明确为 Yes 且许可证属于预设宽松集合来构造操作性“严格开放”样本。原始 4576 条经日期、参数量与六任务筛选及模型去重后，严格开放为 236 个模型，其中预训练 42 个、后训练 141 个；其余类型不混入分层主回归。此筛选可复算，但不等于法律上的开源认证或提交当日已开放证明。
\fig{problem4/figures/q4_sample_attrition.pdf}{0.82}{题目附件 C1/C2 到严格开放样本的筛选过程。}

C1/C4 的算力连接须满足名称、来源身份、参数、发布时间、领域和权重状态的一致性；得到 29 个含算力的预训练连接，其中外部来源可支持的敏感性子集为 15 个、6 个家族。该子集最高算力 $2.94\times10^{23}$ FLOPs，仍低于后文的情景起点。C4 的数据量与权重字段按缺失和口径一致性审计；能与严格开放后训练得分可靠成对连接的训练算力记录为零，故不将 C4 数据量缺失值补成训练 token 数，也不拟合后训练层的联合 $(N,D,C)$ 能力模型。外部来源核查只对连接做补充敏感性，不改写题目附件主字段。C3 的 4573 条排行榜行可匹配 C1；另 26 条早年文献补录的 Average 与六任务均值口径不一致，不能拼为 2019 年以来的同协议长时间序列。C8 的 1954 个 JSON 可做 BBH 叶任务汇总，但汇总分数与 C1 的 BBH 数值标尺未证实相同。以上共同构成样本、跨表和评分的硬边界。

\fig{problem4/figures/q4_observed_frontier.pdf}{0.83}{严格开放模型在 C1 同协议下的观测前沿，按模型类型分层。}

\subsection{规模与时间的条件分解}
将六任务均分映射到 logit 空间，对 $\log_{10}N$ 及排行榜提交时间用 Ridge 拟合；按模型家族分组选择正则化，2025 年 1--3 月作未来时段留出，以家族等权误差评估。严格后训练层中，仅参数量模型家族均衡 RMSE 为 8.0189，加入时间后为 11.736，差值 3.717 的家族 Bootstrap 95\% 区间为 $[-1.113,7.535]$，同时覆盖改善与恶化，故既不能证明时间项更优，也不能在未预设等价界值时宣称两模型等价。严格预训练晚期样本仅 5 个模型，也不足以支撑长期趋势。

在早晚期共同参数支持范围内，对严格后训练层做两因素 Shapley 描述分解：规模项 $-4.876$ 分、条件时间残差 $+5.768$ 分，模型合计 $+0.892$ 分，而样本观测均分变化为 $+2.316$ 分。模型总变化的家族 Bootstrap 区间跨 0；时间项混合架构、数据、后训练和样本选择差异，不是纯技术进步贡献，也不能计算稳定的因果比例。剔除 MATH Lvl 5 后早晚差为 1.173 分，而该单任务差为 14.269 分，显示综合均分对任务构成敏感。
\fig{problem4/figures/q4_conditional_decomposition.pdf}{0.83}{共同参数范围内的描述性规模与时间残差分解。}

\subsection{前沿分位与综合分数的补充检验}
题目关注能力前沿，均分回归不能直接替代条件上分位。为此在严格开放后训练层尝试 $0.9$ 条件分位数回归（它不是最高值），仍以 2025-01 前 74 个模型训练，按 14 个家族分组选择正则化参数，并将 2025-01 至 03 月的 67 个模型（6 家族）整体留出。常数、仅参数量及参数量加时间三类候选的训练期家族均衡 Pinball 损失分别为 2.104、1.803、1.713；训练期选中含时间模型。但在未来时段，其家族均衡损失为 1.322，高于仅参数量的 1.073；逐模型经验覆盖率分别为 0.970 和 0.881，训练分组验证中选中模型的逐模型覆盖率仅 0.689。分位水平为 0.9 不意味着这些异质家族样本具有 90\% 覆盖保证。该补充模型只检验短期前沿口径，不替换正文的描述性分解，也不生成 12/24 个月分数。
\fig{problem4/figures/q4_frontier_quantile_holdout.pdf}{0.86}{条件 90\% 分位候选在未来时段留出上的预测与经验覆盖率；仅有 6 个留出家族。}

六任务均分保留为 C1 的同协议主指标。另以训练期各任务中位数及四分位距固定标准化标尺，早期至晚期的六任务平均变化为 0.202 个训练期四分位距单位；剔除 MATH Lvl 5 后为 0.118。原始单位下六任务均分变化为 3.356 分，二者量纲不同，不直接比较数值大小。MATH Lvl 5 的单任务标准化变化为 0.622，仍提示综合分数对该任务敏感。固定训练期标尺避免用未来样本重设刻度，但没有消除模型家族组成的变化。

在 C1/C4 的 29 个预训练连接上，探索性家族留出比较给出仅算力与参数量加算力模型的 RMSE 分别为 6.691 和 7.117；在来源支持的 15 个连接上分别为 5.178、5.112。结果方向不稳定，且严格开放后训练连接为零，故不能据此分配算力与技术的因果贡献。C6 的中可比 68 点分散在 30 个原始来源标签、其中 14 个标签只有一个点；将它们扩充到高可比 7 点的桥接式会引入未识别的协议偏差，故维持原桥接边界。
\subsection{Loss 桥接及 12/24 个月情景}
C6 与标度律高可比的仅有同一家族 Pythia 的 7 个点。逐点留一的单调 Loss--评分映射 RMSE 为 0.4366，尚差于常数基线 0.4244，且完全没有跨家族验证。因此不能把问题三的 Loss 优化直接换算成排行榜分数。C4 以截至 2025-03-13 的 2025 年不完整年份中位算力 $5.7144\times10^{23}$ FLOPs 为参考，按题给“算力放缓”构造停滞、历史增速减半等条件路径。12 个月的条件算力为 $5.7144\times10^{23}$ 至 $1.0610\times10^{24}$ FLOPs，24 个月为 $5.7144\times10^{23}$ 至 $1.9699\times10^{24}$ FLOPs；历史增速中点仅作对照。改用 2024 年完整年份中位数，所有路径降为上述数值的 0.3323 倍，说明起点选取高度影响情景。

C1/C4 可靠连接的 29 个算力--得分点全部是预训练模型，严格开放后训练的成对点为零。基于这一断裂以及未来留出表现，12/24 个月的具体能力分数在附件下不可识别。可将桥接断点写为 $S_f=h(L)+u_f$，其中 $u_f$ 是家族特有的协议与能力偏移。C6 高可比数据只识别 Pythia 家族内的局部关系，不约束未来新家族的 $u_f$。对任意 $s\in[0,100]$，均可选取一个新家族偏移使 $S_f=s$ 而不改变现有 C6 观测；故在不加跨家族可迁移假设时，新模型分数的识别集为 $[0,100]$。若前沿定义为已有纪录与未来分数之最大值，识别集才缩为 $[46.889,100]$，并非统计置信区间。
图中仅展示算力假设和识别状态，不绘制虚构能力曲线。
\fig{problem4/figures/q4_compute_scenario_identification.pdf}{0.85}{12/24 个月条件算力路径及能力分数不可识别的审计结果。}
